\documentclass[UTF8,a4paper,zihao=-4]{ctexart}
\usepackage{amsmath,amssymb,amsthm}
\usepackage{graphicx,booktabs,tabularx}
\usepackage{geometry}
\geometry{top=2.5cm,bottom=2.5cm,left=3cm,right=2.5cm}
\usepackage{fancyhdr}
\usepackage{caption}
\captionsetup{font=small,labelsep=quad}
\usepackage{listings}
\usepackage[hidelinks]{hyperref}
\linespread{1.3}
\setcounter{secnumdepth}{4}
\setcounter{tocdepth}{3}
\pagestyle{fancy}
\fancyhf{}
\fancyhead[C]{2026 年全国大学生数学建模竞赛}
\fancyfoot[C]{\thepage}
\renewcommand{\headrulewidth}{0.4pt}
\lstset{basicstyle=\ttfamily\small,numbers=left,frame=single,breaklines=true}
\begin{document}
\section{模型建立}
\subsection{目标函数}
\subsubsection{线性化处理}
\paragraph{绝对值拆解} 测试四级编号。
\begin{align}
  z &= 4x_1 + 3x_2 \\
  2x_1 + x_2 &\le 10
\end{align}
\begin{table}[htbp]
\centering
\caption{三种方案的评价结果}
\begin{tabular}{lccc}
\toprule
方案 & 成本 & 质量 & 综合得分 \\
\midrule
S1 & 0.90 & 0.50 & 0.685 \\
S3 & 0.60 & 0.90 & 0.730 \\
\bottomrule
\end{tabular}
\end{table}
\begin{lstlisting}[language=Python]
import numpy as np
print("ok")
\end{lstlisting}
\end{document}
